\chapter{Comment fonctionne la rémunération}\label{ch:in:how-pay-works}

Chaque mois de mars, le contrôleur de l'État de New York estime le bonus moyen versé aux salariés de l'industrie financière de la ville
pendant la saison des bonus qui vient de s'achever. Pour 2021, il s'élevait à 257\,500~\$~; pour 2022, 180\,000~\$, soit 30~\% de moins~; pour 2024, 244\,700~\$,
31.5~\% de plus que l'année précédente. Un bonus moyen n'est qu'un gros titre, et il masque la structure qui décide de ce que vaut
la rémunération pour celui qui la reçoit~: quelle part est fixe, quelle part dépend d'une formule ou du jugement d'un manager,
quelle part est différée et sous quelle forme, et ce qui est perdu au départ. Ce chapitre décompose la rémunération du point de vue du salarié.
Le livre~16, chapitre~10, a conçu la même mécanique du côté de la firme, avec ses enveloppes de bonus, ses formules de reversement et ses calendriers de report~;
ici la question est de savoir ce que vaut une offre, et pour qui.

\section{Fixe et bonus~: la structure en deux parties et sa raison d'être}

\begin{definition}[Salaire de base, rémunération totale]\label{def:in:how-pay-works:base}
Le \emph{salaire de base}\index{salaire de base} est la rémunération annuelle fixe, prévue par le contrat de travail et versée par mensualités
quels que soient les résultats de l'année. La \emph{rémunération totale}\index{rémunération totale} est le salaire de base auquel s'ajoute la part variable
attribuée pour une année, en numéraire et différée, à sa valeur d'attribution, les avantages étant présentés séparément~; c'est le
chiffre sur lequel les offres sont comparées, et il compte les attributions différées avant toute \omterm{def:in:how-pay-works:rsu}{déchéance} ou variation de leur valeur.
\end{definition}

Dans le trading, la rémunération comporte deux parties parce que le revenu de la firme en comporte deux~: une base de coûts qu'elle doit couvrir chaque année, et un
résultat qui fluctue avec les marchés (le chapitre~12 a mesuré cette fluctuation selon le type de firme). Une attribution variable permet à la firme de partager les
bonnes années et de réduire la voilure dans les mauvaises sans réduire ses effectifs, ce qui est le ratio de rémunération du livre~16 à l'œuvre. Pour le
salarié, la même structure transfère le risque~: le fixe constitue un plancher, et le reste évolue avec l'année de la firme, celle du desk
et celle de l'individu.

\begin{dated}{2026-03}{Bonus de l'industrie financière de la ville de New York}\label{dat:in:how-pay-works:nyc}
\footnotesize
\begin{tabular}{@{}lrrrrr@{}}
\toprule
année du bonus & 2021 & 2022 & 2023 & 2024 & 2025\\
\midrule
bonus moyen, première estimation (\$) & 257\,500 & 180\,000 & 176\,500 & 244\,700 & 246\,900\\
\bottomrule
\end{tabular}

\smallskip\noindent Contrôleur de l'État de New York, estimations annuelles publiées chaque mois de mars (2022--2026)~; chiffre propre à chaque publication pour son
année, sauf pour 2022, qui est le chiffre donné pour l'année précédente par la publication de 2024. La publication de 2026
indique que la moyenne 2025 est supérieure de 6~\% à celle de 2024, de sorte que l'estimation 2024 avait été révisée à la baisse. L'enveloppe de bonus a atteint un record
de 49.2~milliards de dollars pour 2025. En 2024, le salaire moyen de l'industrie financière de la ville, bonus compris, s'élevait à 505\,677~\$, et les
bonus représentaient environ 42~\% des salaires de l'industrie.
\end{dated}

Deux éléments de l'encadré comptent pour la suite. La moyenne mêle tous les métiers de l'industrie, depuis les opérations jusqu'aux
traders seniors, si bien qu'elle ne décrit personne~; et les bonus représentaient environ deux cinquièmes des salaires de l'industrie, de sorte que la part variable n'est pas
un détail. Le chapitre~14 donne les niveaux~; le présent chapitre en donne la structure.

\section{Rémunération à la formule et discrétionnaire, vue du salarié}

Une attribution variable est soit calculée par une formule à partir d'un résultat mesuré, soit fixée par le jugement d'un manager dans une enveloppe
(livre~16, chapitre~10, définit la \emph{rémunération à la formule} et le \emph{bonus discrétionnaire}). Les plateformes
du chapitre~5 rémunèrent les gérants de portefeuille à un taux de reversement sur le profit de leur carnet (livre~16, chapitre~3)~; le salarié d'une banque ou
d'un teneur de marché est en général payé à partir d'une enveloppe discrétionnaire. Vue du salarié, la différence tient à trois choses.
\begin{itemize}
\item \textbf{Forme.} Un reversement d'une part du profit positif, et rien en dessous de zéro, est une option sur le résultat du carnet~:
convexe, avec une longue queue droite et une masse en zéro. Une attribution discrétionnaire est plus lisse, car le manager répartit l'enveloppe et
pèse plus d'une année.
\item \textbf{Quel risque.} Une formule expose le salarié au résultat du carnet, chance comprise~; la discrétion expose le
salarié à l'enveloppe, qui suit l'année de la firme et celle du desk, ainsi qu'au jugement du manager.
\item \textbf{Ce qui y met fin.} Un carnet à formule peut être fermé par une limite de perte (les échelles de drawdown du chapitre~5), ce qui met fin à la fois au
poste et au reversement~; l'attribution d'un salarié discrétionnaire peut tomber à zéro une mauvaise année sans que le poste prenne fin.
\end{itemize}
L'équivalent certain (livre~4, chapitre~9) traduit ces différences en argent~: le montant sûr qu'une personne
ayant une aversion relative au risque donnée accepterait à la place du paquet incertain. Deux personnes à qui l'on offre le même paquet peuvent
lui donner une valeur différente, parce que leur aversion au risque et leur autre patrimoine diffèrent.

\section{Report, acquisition, déchéance et récupération~: ce que vaut la rémunération différée}

Une partie de l'attribution variable est souvent différée~: attribuée maintenant, versée plus tard par tranches selon un calendrier d'acquisition, et détenue
dans l'intervalle en numéraire, en actions de la firme ou en parts de ses fonds (livre~16, chapitre~10, définit la rémunération différée,
le calendrier d'acquisition, le malus et la récupération).

\begin{definition}[Action de performance, déchéance, clause de good leaver]\label{def:in:how-pay-works:rsu}
Une \emph{action de performance}\index{action de performance} (restricted stock unit, RSU) est une promesse de livrer une action de l'employeur
ou sa valeur en numéraire, à l'acquisition de l'unité, sous réserve de conditions telles que le maintien dans l'emploi. La \emph{déchéance}
\index{déchéance} est la perte des attributions non acquises, généralement lors d'un départ volontaire de la firme ou d'un licenciement pour faute.
Une \emph{clause de good leaver}\index{clause de good leaver} est une stipulation d'une attribution qui permet à certains partants, typiquement en cas de retraite,
de décès, d'invalidité ou de licenciement économique, de conserver les attributions non acquises, qui s'acquièrent alors selon le calendrier initial ou immédiatement.
\end{definition}

\begin{dated}{2025-10}{Règles de report pour les banques et les attributions d'une banque}\label{dat:in:how-pay-works:rules}
\textbf{UE} (directive 2013/36/UE, article~94, telle que modifiée en 2019)~: pour le personnel dont l'activité a un impact significatif sur le
risque de la firme, la rémunération variable ne peut excéder 100~\% de la rémunération fixe (200~\% avec l'approbation des actionnaires)~; au moins 50~\% de la rémunération
variable est versée en actions ou instruments équivalents~; au moins 40~\% est différée sur une période d'au moins quatre à cinq ans, sans acquisition plus
rapide que le prorata, et au moins 60~\% d'un montant particulièrement élevé. \textbf{Royaume-Uni} (PRA PS21/25 et FCA PS25/15,
en vigueur pour les exercices de performance commençant après le 16~octobre 2025)~: un report minimal de quatre ans pour tous les preneurs de
risques significatifs~; 40~\% de report sur les premières 660\,000~\pounds d'un bonus et 60~\% sur la part au-delà~; le plafonnement des bonus a été supprimé à compter du 31~octobre 2023.
\textbf{Une banque américaine} (son formulaire 10-K pour 2025)~: les RSU sans conditions de performance «~sont généralement acquises
et les actions ordinaires sous-jacentes sont livrées (nettes de la retenue à la source requise) sur une période de trois ans~»,
avec une acquisition accélérée en cas de «~retraite, décès, invalidité et, dans certains cas, emploi en conflit d'intérêts~».
\end{dated}

Ce que vaut une attribution différée pour le salarié dépend de quatre éléments, que le modèle du chapitre évalue ensemble.
\begin{enumerate}
\item \textbf{Le temps.} Un dollar dans quatre ans vaut moins qu'un dollar aujourd'hui~; à 5~\% par an, une tranche acquise dans quatre
ans vaut 82~cents d'attribution.
\item \textbf{L'instrument.} Les actions différées portent le risque du titre~: avec une volatilité de 30~\%, la valeur d'une tranche à
l'acquisition varie fortement, même sans variation espérée.
\item \textbf{Le départ.} Les tranches non acquises sont perdues en cas de départ volontaire, si bien que leur valeur est pondérée par la
probabilité de rester~; une \omterm{def:in:how-pay-works:rsu}{clause de good leaver} supprime cette pondération pour les partants qu'elle couvre.
\item \textbf{Malus et récupération.} Une attribution peut être réduite avant son acquisition (malus) ou récupérée après celle-ci (clawback) en cas de
faute ou de perte ultérieure~; le modèle du chapitre traite ces cas comme rares et les énonce sans les chiffrer.
\end{enumerate}
L'impôt suit l'acquisition, et non l'attribution. Aux États-Unis, un bien soumis à «~un risque substantiel de \omterm{def:in:how-pay-works:rsu}{déchéance}~»
n'est pas un revenu tant qu'il «~n'est pas devenu substantiellement acquis~», et c'est alors sa valeur de marché qui devient imposable~; l'impôt sur une attribution différée est payé
sur sa valeur à l'acquisition, ce qui explique qu'une banque livre des actions «~nettes de la retenue à la source requise~». Le chapitre~15 expose les
règles des autres pays.

\begin{example}[La menotte]\label{ex:in:how-pay-works:handcuff}
L'offre bancaire illustrative du chapitre diffère 40~\% de chaque attribution en actions de la firme, acquises par tranches égales sur
quatre ans. Pour celui qui reste, le solde non acquis, à la valeur d'attribution, atteint environ 90\,000~\$ à la fin de la première
année, 158\,000~\$ après deux ans, 203\,000~\$ après trois ans, puis se stabilise autour de 226\,000~\$ à partir de la quatrième année, moment où le nouveau
report de chaque année est compensé par les acquisitions (\cref{fig:in:how-pay-works:unvested}). Un départ à la fin de la deuxième année fait perdre une valeur actuelle
d'environ 129\,000~\$~; à la fin de la quatrième année, 170\,000~\$.
\end{example}

\begin{omfigure}
\begin{tikzpicture}
\begin{axis}[width=10.5cm,height=5.4cm,ybar,area legend,bar width=10pt,xmin=0.4,xmax=5.6,ymin=0,ymax=260,xtick={1,2,3,4,5},
  xlabel={\small fin d'année},ylabel={\small milliers de \$},tick label style={font=\footnotesize},grid=major,
  grid style={gray!20},table/col sep=comma,legend style={legend cell align=left,font=\scriptsize,at={(0.5,-0.3)},anchor=north,
  /tikz/every even column/.append style={column sep=8pt}},legend columns=2]
\addplot[fill=omDef!55,draw=omDef] table[x=year,y=unvested]{figdata/industry/13-how-pay-works/bank_unvested.csv};
\addplot[fill=omThm!55,draw=omThm] table[x=year,y=forfeit]{figdata/industry/13-how-pay-works/bank_unvested.csv};
\legend{solde non acquis à la valeur d'attribution (rester),valeur actuelle perdue en partant alors}
\end{axis}
\end{tikzpicture}

\omcaption{Les attributions différées de l'offre bancaire illustrative~: 40~\% de l'attribution annuelle est différée en actions et
acquises sur quatre ans. Le solde qu'un partant abandonne se constitue pendant quatre ans, puis se stabilise. Les paramètres sont
illustratifs~; données~: 20\,000 carrières simulées, via \texttt{in\_pay.bank\_unvested} et \texttt{forfeit\_by\_year}.}\label{fig:in:how-pay-works:unvested}
\end{omfigure}

\section{Primes d'arrivée, garanties et rachats}

\begin{definition}[Prime d'arrivée, bonus garanti]\label{def:in:how-pay-works:signon}
Une \emph{prime d'arrivée}\index{prime d'arrivée} (sign-on bonus) est un versement à l'entrée, généralement remboursable en totalité ou en partie si le
salarié quitte la firme dans un délai fixé. Un \emph{bonus garanti}\index{bonus garanti} (guaranteed bonus) est une rémunération variable dont le
minimum pour une année donnée, en général la première, est fixé dans l'offre.
\end{definition}

Une firme qui recrute utilise les deux instruments pour résoudre deux problèmes. La \omterm{def:in:how-pay-works:signon}{prime d'arrivée} compense ce que le candidat abandonne en changeant
de poste, et sa clause de remboursement le retient pendant sa durée~; la garantie écarte le risque de la première année, celle où un nouvel
arrivant n'a pas encore d'historique dans la firme et ne dispose que d'une année partielle pour en constituer un. Un troisième instrument, le \emph{rachat des différés} (buyout) du livre~16,
chapitre~10, remplace les attributions non acquises que le candidat perd en partant, en général sous forme de nouvelles attributions aux anciennes
dates d'acquisition. Du côté du candidat, le calcul est la menotte de l'exemple ci-dessus lue à l'envers~: un candidat qui
quitte la banque illustrative à la fin de la deuxième année perd des attributions d'une valeur d'attribution d'environ 158\,000~\$, et un rachat
de cette taille aux mêmes dates rend le changement neutre pour la partie différée. Le candidat doit aussi demander quelles
conditions de remboursement et de \omterm{def:in:how-pay-works:rsu}{déchéance} s'appliquent au rachat lui-même, car un rachat est en général à son tour différé.

Les règles européennes pour les banques limitent ce qu'une garantie peut faire~: la rémunération variable garantie «~est exceptionnelle,
n'intervient que lors de l'embauche de nouveaux salariés, lorsque l'établissement dispose d'une base de capital solide et forte, et se
limite à la première année d'emploi~»~; en dehors des firmes que ces règles couvrent, elle relève du contrat. Le chapitre~28 revient sur
les rachats et les préavis lorsqu'il suit des changements de firme, et le livre~18 sur la manière dont les offres se négocient.

\section{Associés, actions et carried interest}

\begin{definition}[Part d'associé, carried interest]\label{def:in:how-pay-works:carry}
Une \emph{part d'associé}\index{part d'associé} est le droit d'un associé à une fraction du bénéfice d'une société de personnes, fixée par
le pacte d'associés, en contrepartie du capital apporté et du travail~; elle est versée comme \omterm{def:in:reading-a-trading-firms-accounts:members}{rémunération des associés}
(chapitre~11) plutôt que comme salaire. Le \emph{carried interest}\index{carried interest} est une fraction des bénéfices d'un fonds,
généralement au-delà d'un seuil de rendement, versée à ses gérants par une participation dans le commandité du fonds (general partner) plutôt que sous forme de commission.
\end{definition}

Être propriétaire rémunère autrement que d'être salarié. Un associé d'une société de personnes à responsabilité limitée (livre~16, chapitre~2)
partage le bénéfice après les coûts, y compris les salaires des employés, et supporte ses pertes à hauteur du capital apporté~; le chapitre~11
a montré que, dans certaines sociétés de personnes de trading, l'essentiel du bénéfice revient à un associé personne morale plutôt qu'aux individus.
Les actions d'une firme privée valent ce que la firme acceptera d'en payer selon les termes de ses accords, qui restreignent en général la
vente et fixent un prix au départ. Le \omterm{def:in:how-pay-works:carry}{carried interest} est la part d'une commission de performance versée aux personnes qui gèrent
le fonds (livre~1, chapitre~3, définit la commission et le high-water mark).

\begin{dated}{2026-04}{Taxation du carried interest}\label{dat:in:how-pay-works:carry}
\textbf{Royaume-Uni}~: à partir du 6 avril 2026, le \omterm{def:in:how-pay-works:carry}{carried interest} est imposé entièrement dans l'impôt sur le revenu, «~traité comme des bénéfices
de trading et soumis à l'Income Tax et aux cotisations de Class 4 National Insurance~»~; pour le \omterm{def:in:how-pay-works:carry}{carried interest} éligible, «~la part à traiter
comme bénéfices de trading est de 72.5~\% des bénéfices éligibles~». \textbf{États-Unis}~: les gains sur un intérêt de société de personnes
applicable ne sont à long terme qu'après trois ans de détention au lieu d'un (26 U.S.C.\ 1061).
\end{dated}

\section{Structure selon le type de firme}

La même combinaison d'éléments varie selon le type d'employeur, pour les raisons exposées dans les chapitres précédents.
\begin{itemize}
\item \textbf{Banques.} La rémunération variable au-delà des seuils réglementaires est différée et versée en partie en actions, selon les règles
de l'encadré daté pour le personnel qu'elles couvrent~; les RSU d'une banque américaine s'acquièrent sur trois ans. Le différé retient, et un
rachat fait partie du coût de l'embauche d'un banquier expérimenté (divisions du chapitre~7, rôles du chapitre~18).
\item \textbf{Teneurs de marché et firmes pour compte propre.} Les firmes privées ne publient aucune règle de rémunération~; une offre comprend en général
un fixe et une attribution discrétionnaire tirée d'une enveloppe qui suit le chiffre d'affaires de la firme, lequel, selon le chapitre~12, varie avec
la volatilité plus que pour tout autre type de firme. Tout différé est contractuel, non réglementaire, en dehors des entités que les règles couvrent.
\item \textbf{Plateformes.} La rémunération d'un gérant de portefeuille est un taux de reversement sur le bénéfice du book (chapitre~5)~; celle d'un
analyste, dans l'offre illustrative de ce chapitre, est un fixe et une petite part du résultat du book~; le book peut être fermé par une
limite de perte. Les \omterm{def:in:how-pay-works:signon}{primes d'arrivée} et les garanties rachètent les personnes à leur précédente firme.
\item \textbf{Sociétés de gestion et fonds.} Fixe et bonus, qui peuvent être en partie différés dans les fonds ou les actions de la firme~;
les cadres supérieurs peuvent partager les commissions de performance ou le \omterm{def:in:how-pay-works:carry}{carried interest} (chapitre~8).
\item \textbf{Bourses, fournisseurs de données et régulateurs.} Les bourses cotées et les fournisseurs versent un fixe et un bonus comme les autres
sociétés cotées~; un régulateur rémunère selon une grille publiée (chapitre~9).
\item \textbf{Firmes crypto.} Les jetons et les \omterm{def:in:crypto-native-firms:tokens}{bons de souscription de jetons} (token warrant) ajoutent une composante volatile, souvent bloquée (chapitre~10).
\end{itemize}

\section{Tutoriel~: trois offres, un choix}\label{tut:in:how-pay-works:tut}

\begin{tutorial}
\textbf{Objectif.} Valoriser trois offres stylisées du point de vue du salarié et les comparer à deux niveaux d'aversion au risque.
\textbf{État final~:} \cref{fig:in:how-pay-works:offers,fig:in:how-pay-works:ce} et le tableau ci-dessous.
\begin{tutsteps}
\item \textbf{Les offres} (tous les nombres sont illustratifs, en dollars américains)~: un poste de \emph{banque} avec un fixe de 250\,000 et une
attribution discrétionnaire de médiane 200\,000 (écart-type logarithmique 0.5), 40~\% différés en actions (volatilité 30~\%) sur quatre
ans, déchus en cas de départ~; un poste de \emph{teneur de marché} (market maker) avec un fixe de 200\,000 et une attribution de médiane 250\,000 (écart-type logarithmique
0.8), versée en espèces~; un \emph{analyste de plateforme} avec un fixe de 175\,000, un versement de 2~\% du bénéfice du book (moyenne
15 millions, écart-type 20 millions par an), le book fermé après une année de perte supérieure à 10 millions, une \omterm{def:in:how-pay-works:signon}{prime d'arrivée} de
100\,000 remboursable dans l'année, et une garantie de la première année de 150\,000.
\item \textbf{La carrière.} Cinq ans, une probabilité annuelle de 10~\% de quitter volontairement l'entreprise, un taux d'actualisation de 5~\%~;
\texttt{firm.payoffer.simulate} tire 20\,000 carrières par offre (\cref{lst:in:how-pay-works:sim}).
\omcode{code/firm/payoffer/firm_payoffer.py}{87}{103}{Chaque année~: le fixe pendant le travail, la part versée d'avance de l'attribution,
et chaque tranche différée livrée si le salarié est encore en poste ou s'il est good leaver, déchue sinon.}\label{lst:in:how-pay-works:sim}
\item \textbf{La valeur.} \texttt{summary} donne la moyenne et les centiles de la valeur actuelle reçue~;
\texttt{certainty\_equivalent(pv, rra, 500\_000)} donne l'équivalent certain avec 500\,000 d'autres richesses.
\end{tutsteps}
\end{tutorial}

\begin{center}\footnotesize
\begin{tabular}{@{}lrrrrr@{}}
\toprule
offre (illustrative) & VA moyenne & 10\up{e}~centile & 90\up{e}~centile & CE, RRA 1 & CE, RRA 3\\
\midrule
banque & 1\,618 & 597 & 2\,230 & 1\,500 & 1\,186\\
teneur de marché & 1\,949 & 686 & 2\,956 & 1\,779 & 1\,365\\
analyste de plateforme & 1\,703 & 171 & 3\,010 & 1\,439 & 904\\
\bottomrule
\end{tabular}

\smallskip Valeur actuelle de cinq ans de rémunération, en milliers de dollars~; CE est l'équivalent certain.
\end{center}

\begin{omfigure}
\begin{tikzpicture}
\begin{axis}[width=10.5cm,height=5.4cm,xmin=0.4,xmax=3.6,ymin=0,ymax=3400,xtick={1,2,3},
  xticklabels={banque,{teneur de marché},{analyste de\\plateforme}},xticklabel style={align=center},ylabel={\small valeur actuelle, en milliers de dollars},
  tick label style={font=\footnotesize},grid=major,grid style={gray!20},table/col sep=comma,
  legend style={legend cell align=left,font=\scriptsize,at={(0.5,-0.2)},anchor=north,/tikz/every even column/.append style={column sep=8pt}},
  legend columns=2]
\addplot[draw=none,mark=none,forget plot,error bars/.cd,y dir=both,y explicit,error mark=none,
  error bar style={omDef,line width=1pt}]
  table[x=pos,y=p50,y error plus expr=\thisrow{p90}-\thisrow{p50},y error minus expr=\thisrow{p50}-\thisrow{p10}]{figdata/industry/13-how-pay-works/dist.csv};
\addplot[draw=none,mark=none,forget plot,error bars/.cd,y dir=both,y explicit,error mark=none,
  error bar style={omDef!35,line width=14pt}]
  table[x=pos,y=p50,y error plus expr=\thisrow{p75}-\thisrow{p50},y error minus expr=\thisrow{p50}-\thisrow{p25}]{figdata/industry/13-how-pay-works/dist.csv};
\addplot[only marks,mark=-,mark size=7pt,line width=1.4pt,omDef] table[x=pos,y=p50]{figdata/industry/13-how-pay-works/dist.csv};
\addplot[only marks,mark=diamond*,mark size=2.5pt,omThm] table[x=pos,y=mean]{figdata/industry/13-how-pay-works/dist.csv};
\legend{médiane,moyenne}
\end{axis}
\end{tikzpicture}

\omcaption{Cinq ans de rémunération sous trois offres illustratives~: les boîtes couvrent les 25\up{e} à 75\up{e}~centiles de la valeur actuelle
perçue, les moustaches les 10\up{e} à 90\up{e}~centiles, la barre la médiane et le losange la moyenne. Le spread de l'analyste de plateforme est le plus large~: un livre clôturé
met fin à la fonction et au reversement en même temps. Les paramètres sont illustratifs~; données~: \texttt{in\_pay.table}.}\label{fig:in:how-pay-works:offers}
\end{omfigure}

\begin{omfigure}
\begin{tikzpicture}
\begin{axis}[width=10.5cm,height=5.4cm,xmin=0.4,xmax=5.1,ymin=500,ymax=2000,xlabel={\small aversion relative au risque},
  ylabel={\small équivalent certain, en milliers de dollars},tick label style={font=\footnotesize},grid=major,grid style={gray!20},
  table/col sep=comma,legend style={legend cell align=left,font=\scriptsize,at={(0.97,0.97)},anchor=north east}]
\addplot[omDef,thick,mark=*] table[x=rra,y=bank]{figdata/industry/13-how-pay-works/ce_curve.csv};
\addplot[omThm,thick,mark=square*] table[x=rra,y=maker]{figdata/industry/13-how-pay-works/ce_curve.csv};
\addplot[black!70,thick,mark=triangle*] table[x=rra,y=platform]{figdata/industry/13-how-pay-works/ce_curve.csv};
\legend{banque,teneur de marché,analyste de plateforme}
\end{axis}
\end{tikzpicture}

\omcaption{Équivalents certains des trois offres illustratives en fonction de l'aversion relative au risque, avec 500\,000 de
patrimoine par ailleurs. Le package de plateforme perd le plus à mesure que l'aversion au risque augmente~: à faible aversion, il vaut à peu près autant que l'offre
de la banque, à forte aversion, beaucoup moins. Données~: \texttt{in\_pay.ce\_curve}.}\label{fig:in:how-pay-works:ce}
\end{omfigure}

L'offre du teneur de marché vaut le plus à chaque niveau d'aversion au risque retenu ici, parce que ses paramètres lui donnent
la prime médiane la plus élevée~; c'est une propriété des chiffres illustratifs, non des teneurs de marché. La leçon est dans les formes. La
moyenne du package de plateforme est supérieure à celle de la banque, mais son 10\up{e}~centile est de 171\,000~\$, car un livre clôturé après une mauvaise année
met fin à la rémunération~; dans 35.4~\% des carrières simulées, le livre est clôturé dans les cinq ans. Son équivalent certain tombe de
1\,439\,000~\$ à 904\,000~\$ quand l'aversion relative au risque passe de~1 à~3, tandis que celui de la banque passe de 1\,500\,000~\$ à 1\,186\,000~\$.
Un candidat qui ne peut pas supporter la queue basse de la plateforme doit l'évaluer nettement en dessous de sa moyenne.

\textbf{Pour aller plus loin.} Donner à l'offre de la banque une \omterm{def:in:how-pay-works:rsu}{clause de good leaver} et mesurer la part de la \omterm{def:in:how-pay-works:rsu}{déchéance} qu'elle supprime
pour un départ prévu~; relever la limite de perte de la plateforme et suivre l'équivalent certain.

\section{Construction~: l'évaluateur d'offres}\label{bld:in:how-pay-works:payoffer}

\begin{build}
\bfield{Rôle} Évaluer un package de rémunération du point de vue du salarié, pour ce chapitre et pour les chapitres~22, 28 et~30 (et pour
le chapitre du livre~18 consacré aux offres).
\bfield{Interface} \texttt{firm.payoffer}~:
\texttt{Package(name, base, median, sigma, deferral, vest\_years, inst\_vol, good\_leaver, sign\_on, sign\_on\_years, guarantee, share, pnl\_mean, pnl\_sd, cut)}~;
\texttt{simulate(pkg, years, hazard, rate, n, rng, leave\_year)}~; \texttt{summary(sim)}~;
\texttt{certainty\_equivalent(pv, rra, wealth)}~; \texttt{buyout(pkg, year, n, rng, rate)}.
\bfield{Règles} Les primes sont versées à la fin de l'année travaillée~; les tranches différées sont acquises à parts égales après la prime et
ne sont livrées qu'aux personnes encore en poste ou couvertes comme good leavers~; une \omterm{def:in:how-pay-works:signon}{prime d'arrivée} est remboursée intégralement sur sa période~;
un \omterm{def:in:how-pay-works:signon}{bonus garanti} fixe un plancher à la première prime~; un livre clôturé ne verse rien pour sa dernière année~; chaque paramètre relève de l'appelant.
\bfield{Tests d'acceptation} \texttt{code/firm/payoffer/tests/}~: un package de base seul est une annuité~; un différé sans risque
correspond au calcul à la main~; un départ fait perdre exactement les tranches non acquises et un good leaver n'en perd aucune~; une \omterm{def:in:how-pay-works:signon}{prime d'arrivée} est remboursée lors
d'une sortie anticipée~; le rachat égale le solde non acquis~; l'équivalent certain d'un montant sûr est ce montant et diminue
avec l'aversion au risque~; la moyenne avec risque sur action correspond à la forme fermée à 1~\% près.
\bfield{Pour aller plus loin} Malus et récupération (clawback) comme événements de probabilité donnée~; fiscalité selon le lieu (chapitre~15)~; primes corrélées
d'une année à l'autre~; un taux de départ qui dépend de la prime de l'année.
\end{build}

\begin{omsources}
\item Office of the New York State Comptroller, securities industry bonus releases, mars 2022, 2024, 2025 et 2026.
\item Directive 2013/36/EU, article~94, et Directive (EU) 2019/878~; Bank of England, PS21/25 -- Remuneration Reform.
\item The Goldman Sachs Group, Inc., Form 10-K for 2025, note~29.
\item US Internal Revenue Service, Publication 525~; 26 U.S.C.\ 1061.
\item HM Treasury et HMRC, Reform of the tax treatment of carried interest~: revised tax regime (2025).
\end{omsources}

\section{Exercices}

\begin{exercise}[$\star$]\label{exo:in:how-pay-works:1}
De combien le bonus moyen du secteur des valeurs mobilières de la ville de New York a-t-il varié entre 2021 et 2022, puis entre 2023 et 2024~?
\end{exercise}

\begin{exercise}[$\star$]\label{exo:in:how-pay-works:2}
Selon les règles britanniques en vigueur pour les années de performance postérieures à octobre 2025, quelle part d'un bonus de 1\,000\,000~£ doit
être différée~?
\end{exercise}

\begin{exercise}[$\star$]\label{exo:in:how-pay-works:3}
À un taux d'actualisation de 5~\%, que vaut une tranche acquise dans quatre ans par unité de prime, avant tout risque~?
\end{exercise}

\begin{exercise}[$\star\star$]\label{exo:in:how-pay-works:4}
Dans l'offre de la banque, lister les tranches encore non acquises à la fin de la deuxième année et calculer leur valeur d'attribution, en utilisant une
prime attendue de $200\,000\,e^{0.5^2/2}$.
\end{exercise}

\begin{exercise}[$\star\star$]\label{exo:in:how-pay-works:5}
Pourquoi un reversement égal à une part du profit positif est-il une option, et qu'est-ce que cela change à son équivalent certain~?
\end{exercise}

\begin{exercise}[$\star\star$]\label{exo:in:how-pay-works:6}
Une candidate se voit proposer une \omterm{def:in:how-pay-works:signon}{prime d'arrivée} de 100\,000~\$, remboursable intégralement si elle quitte le poste dans l'année. Elle estime à 15~\% la
probabilité qu'elle le fasse. Que vaut la \omterm{def:in:how-pay-works:signon}{prime d'arrivée} pour elle en espérance, sans tenir compte de l'actualisation~?
\end{exercise}

\begin{exercise}[$\star\star\star$]\label{exo:in:how-pay-works:7}
\emph{Programmation.} Donner à l'offre de la banque une \omterm{def:in:how-pay-works:rsu}{clause de good leaver} couvrant les départs volontaires. De combien baisse la
\omterm{def:in:how-pay-works:rsu}{déchéance} espérée, et de combien monte la valeur actuelle moyenne~?
\end{exercise}

\begin{exercise}[$\star\star\star$]\label{exo:in:how-pay-works:8}
\emph{Trouver la faille.} «~Le bonus moyen à Wall Street s'élevait à 246\,900~\$ en 2025, donc un trader rejoignant une banque à New York
doit s'attendre à peu près à ce montant.~»
\end{exercise}

\section{Problème~: trois offres, un choix}

\begin{problem}[{Problème du week-end --- trois offres, un choix}]\label{pb:in:how-pay-works:1}
Une quant disposant de 500\,000 d'épargne détient trois offres~: une banque, un teneur de marché et une plateforme. Elle veut savoir ce que
chacune vaut pour elle, et ce que coûterait à la banque son départ au bout de deux ans.

\medskip\noindent\textbf{Partie I --- Les composantes.}
\begin{enumerate}
\item Définir le \omterm{def:in:how-pay-works:base}{salaire de base} et la \omterm{def:in:how-pay-works:base}{rémunération totale}. Pourquoi les offres se comparent-elles sur la \omterm{def:in:how-pay-works:base}{rémunération totale}, et que laisse-t-elle de côté~?
\item Donner le bonus moyen dans la finance new-yorkaise pour 2021--2025 et la part des salaires que représentaient les bonus en 2024.
\item Pourquoi la rémunération en trading comporte-t-elle deux parties~?
\item Opposez la rémunération à la formule et la rémunération discrétionnaire quant à leur forme, au risque qu'elles portent et à ce qui y met fin.
\item Définir l'équivalent certain et indiquer de quoi il dépend, outre le package.
\end{enumerate}

\medskip\noindent\textbf{Partie II --- Différé.}
\begin{enumerate}[resume]
\item Définir une \omterm{def:in:how-pay-works:rsu}{action de performance} (RSU), la \omterm{def:in:how-pay-works:rsu}{déchéance} et une \omterm{def:in:how-pay-works:rsu}{clause de good leaver}.
\item Énoncer les règles de différé de l'UE et du Royaume-Uni, ainsi que le calendrier d'acquisition des RSU d'une banque américaine.
\item Citer quatre éléments qui fixent la valeur d'une attribution différée pour le salarié.
\item Quand une RSU d'un salarié américain est-elle imposée, et sur quelle valeur~?
\item Donner le solde non acquis de l'offre de la banque à la fin de chaque année, ainsi que la valeur perdue par un départ après la deuxième année.
\end{enumerate}

\medskip\noindent\textbf{Partie III --- Les offres.}
\begin{enumerate}[resume]
\item Définir une \omterm{def:in:how-pay-works:signon}{prime d'arrivée} et un \omterm{def:in:how-pay-works:signon}{bonus garanti}, et indiquer ce que chacun résout pour la société.
\item Énoncer les paramètres des trois offres et les hypothèses de carrière.
\item Donner la valeur actuelle moyenne de chaque offre, ainsi que ses 10\up{e}~centile et 90\up{e}~centile.
\item Dans quelle part des carrières le livre de la plateforme est-il clos au bout de cinq ans~?
\item Définir le \omterm{def:in:how-pay-works:carry}{carried interest} et une \omterm{def:in:how-pay-works:carry}{part d'associé}, et énoncez les règles fiscales du \omterm{def:in:how-pay-works:carry}{carried interest} au Royaume-Uni et aux États-Unis.
\end{enumerate}

\medskip\noindent\textbf{Partie IV --- Le verdict.}
\begin{enumerate}[resume]
\item Énoncer le \emph{résultat nommé}~: l'équivalent certain de chaque offre pour une aversion relative au risque de 1 et de 3, et le
rachat des différés qui rend neutre le départ de la banque la deuxième année.
\item Pourquoi l'offre de la plateforme perd-elle le plus à mesure que l'aversion au risque augmente~?
\item Que doit-elle demander à propos d'un rachat, outre son montant~?
\item Quel résultat dépend des paramètres illustratifs, et lequel dépend de la structure~?
\item En deux phrases, comment doit-elle choisir~?
\end{enumerate}
\end{problem}

\section{Questions d'entretien}

\begin{interviewq}[$\star$ \iqroles{trader, chercheur}]\label{iq:in:how-pay-works:1}
Votre bonus représente 2~\% du P\&L positif de votre carnet. Quelle option est-ce, et sur quoi porte-t-elle~?
\end{interviewq}

\begin{interviewq}[$\star$ \iqroles{banque}]\label{iq:in:how-pay-works:2}
Pourquoi les régulateurs bancaires exigent-ils qu'une partie de la rémunération variable soit différée et versée en actions~?
\end{interviewq}

\begin{interviewq}[$\star\star$ \iqroles{chercheur, risque}]\label{iq:in:how-pay-works:3}
Calculez l'équivalent certain, avec une utilité logarithmique et une richesse de 100, d'un pile ou face qui rapporte 0 ou 200.
\end{interviewq}

\begin{interviewq}[$\star\star$ \iqroles{développeur}]\label{iq:in:how-pay-works:4}
Concevez le modèle de données des attributions d'un salarié de sorte que le solde non acquis à une date quelconque, et la
\omterm{def:in:how-pay-works:rsu}{déchéance} à une date de départ quelconque, puissent être calculés.
\end{interviewq}

\begin{interviewq}[$\star\star$ \iqroles{banque, risque}]\label{iq:in:how-pay-works:5}
Quelle est la différence entre malus et clawback, et lequel est le plus difficile à faire appliquer~?
\end{interviewq}

\begin{interviewq}[$\star\star\star$ \iqroles{chercheur}]\label{iq:in:how-pay-works:6}
Une société diffère 40~\% de la rémunération sur quatre ans. Estimez combien cela lui fait économiser en rotation du personnel, et quelles données il vous faudrait~?
\end{interviewq}
